% RESUMO
\begin{resumo}
Esta pesquisa investiga a integração de práticas de DevSecOps com Infraestrutura como Código (IaC), com o objetivo de propor e validar um modelo automatizado que promova a segurança e a conformidade contínua em ambientes de computação em nuvem. A problemática central fundamenta-se na crescente complexidade das infraestruturas modernas e na necessidade crítica de mitigar vulnerabilidades comuns, como configurações inseguras e vazamento de credenciais, que elevam significativamente os riscos e custos de violações de dados. Por meio de uma abordagem de segurança \textit{Shift-Left}, o trabalho propõe o deslocamento dos controles de segurança para as fases iniciais do ciclo de vida de desenvolvimento, transformando políticas de conformidade em código auditável e executável.

A metodologia caracteriza-se como pesquisa aplicada, de natureza descritiva e exploratória, com abordagem mista (qualitativa e quantitativa), operacionalizada como um estudo de caso técnico com \textbf{duas Provas de Conceito (PoC) complementares}: (i) uma em nuvem pública real na Amazon Web Services (AWS), com VPC segmentada, EKS, RDS multi-AZ, KMS, S3 e federação de identidade via OpenID Connect (OIDC); e (ii) uma em ambiente local baseado em contêineres Docker, que espelha a mesma topologia lógica para garantir reprodutibilidade acadêmica em uma única máquina, sem dependência de conta em provedor de nuvem. O ecossistema tecnológico utilizou Terraform para o provisionamento declarativo e uma esteira de Integração Contínua (CI/CD) em GitHub Actions integrando ferramentas líderes de mercado: Gitleaks e GitGuardian para detecção de segredos em duas camadas (OSS + SaaS), Checkov para análise estática de segurança (SAST) baseada nos \textit{CIS Benchmarks}, Trivy e Hadolint para segurança de imagens e \textit{Dockerfiles}, e \textit{Open Policy Agent} (OPA) para aplicação de regras de governança, conformidade e FinOps na linguagem Rego.

Os resultados obtidos demonstraram a alta eficácia do modelo proposto, que identificou com sucesso \textbf{100\%} dos casos deliberadamente inseguros submetidos aos testes controlados -- por exemplo, exposição de portas SSH (22) e RDP (3389) para a Internet, ausência de criptografia em repouso em \textit{buckets} S3 e em RDS/EBS, contêineres com \texttt{privileged=true} e imagens com \textit{tag} \texttt{:latest} -- e interceptou-os antes de qualquer \texttt{terraform apply}. A análise comparativa revelou um ganho expressivo de agilidade operacional, reduzindo o tempo médio de revisão de segurança de 15--20 minutos (processo manual) para \textbf{menos de um minuto na PoC de contêineres} e poucos minutos na PoC AWS, o que representa uma ordem de grandeza de economia de aproximadamente 20$\times$. A taxa de falsos positivos observada foi baixa (próxima a 5\%), concentrada em ferramentas de detecção de segredos e mitigada por configuração fina em \texttt{.gitleaks.toml} e supressões documentadas em linha. Adicionalmente, a implementação de \textit{Pre-commit Hooks} e de mecanismos \textit{Fail-Fast} bloqueou credenciais expostas antes do \textit{push} ao repositório em aproximadamente 95\% dos casos, evitando processos corretivos custosos como reescrita do histórico do Git.

Conclui-se que a convergência entre DevSecOps e IaC não apenas fortalece a postura de segurança da organização, mas também democratiza a responsabilidade pela segurança entre as equipes de desenvolvimento e operações. O fato de o mesmo modelo conceitual ser expresso em dois substratos radicalmente diferentes -- nuvem pública gerenciada e contêineres locais -- reforça sua generalidade e amplia sua utilidade acadêmica e didática, permitindo que qualquer pesquisador reproduza integralmente os experimentos. O estudo contribui para a área de Engenharia de Software ao fornecer diretrizes práticas e duas implementações de referência publicadas em código aberto para a construção de infraestruturas resilientes, auditáveis e em conformidade com as melhores práticas de governança e proteção de dados.

 \textbf{Palavras-chave}: DevSecOps. Infraestrutura como Código. Terraform. \textit{Policy-as-Code}. \textit{Shift-Left}. Segurança em Nuvem. Contêineres.
\end{resumo}

\newpage


\begin{resumo}[Abstract]
 \begin{otherlanguage*}{english}
  This research investigates the integration of DevSecOps practices with Infrastructure as Code (IaC), aiming to propose and validate an automated model that promotes security and continuous compliance in cloud computing environments. The central problem is based on the growing complexity of modern infrastructures and the critical need to mitigate common vulnerabilities, such as insecure configurations and credential leakage, which significantly increase the risks and costs of data breaches. Through a \textit{Shift-Left} security approach, the study proposes moving security controls to the early stages of the development lifecycle, transforming compliance policies into auditable and executable code.

The methodology is characterized as applied, descriptive-exploratory research with a mixed-method approach, operationalized as a technical case study through \textbf{two complementary Proofs of Concept (PoC)}: (i) one deployed on real Amazon Web Services (AWS) public cloud, featuring a segmented VPC, EKS, multi-AZ RDS, KMS, S3, and identity federation via OpenID Connect (OIDC); and (ii) another running locally on Docker containers, mirroring the same logical topology to ensure academic reproducibility on a single machine, without dependency on any cloud provider account. The technological ecosystem used Terraform for declarative provisioning and a GitHub Actions Continuous Integration (CI/CD) pipeline integrating market-leading tools: Gitleaks and GitGuardian for two-layer (OSS + SaaS) secret detection, Checkov for static application security testing (SAST) based on the \textit{CIS Benchmarks} framework, Trivy and Hadolint for image and \textit{Dockerfile} security, and \textit{Open Policy Agent} (OPA) for enforcing governance, compliance, and FinOps rules through the Rego language.

The results demonstrated the high effectiveness of the proposed model, which successfully identified \textbf{100\%} of the deliberately insecure cases submitted to controlled tests -- e.g. SSH (22) and RDP (3389) ports exposed to the Internet, lack of encryption at rest in S3 buckets and in RDS/EBS, \texttt{privileged=true} containers, and images tagged \texttt{:latest} -- intercepting them prior to any \texttt{terraform apply}. Comparative analysis revealed a significant gain in operational agility, reducing the average security review time from 15--20 minutes (manual process) to \textbf{less than one minute on the container-based PoC} and a few minutes on the AWS PoC, representing an order-of-magnitude improvement of approximately 20$\times$. The false-positive rate was low (around 5\%), concentrated in secret-detection tools and mitigated through fine-tuning in \texttt{.gitleaks.toml} and inline documented suppressions. Additionally, the implementation of \textit{Pre-commit Hooks} and \textit{Fail-Fast} mechanisms blocked exposed credentials before any \textit{push} in roughly 95\% of the cases, avoiding costly corrective procedures such as rewriting Git history.

The study concludes that the convergence between DevSecOps and IaC not only strengthens the organization's security posture but also democratizes security responsibility across development and operations teams. Expressing the same conceptual model on two radically different substrates -- managed public cloud and local containers -- reinforces its generality and broadens its academic and pedagogical value, allowing any researcher to fully reproduce the experiments. This research contributes to the Software Engineering field by providing practical guidelines and two open-source reference implementations for building resilient, auditable infrastructures in compliance with governance and data protection best practices.

  \vspace{\onelineskip}
  \noindent 
  \textbf{Keywords}: DevSecOps. Infrastructure as Code. Terraform. \textit{Policy-as-Code}. \textit{Shift-Left}. Cloud Security. Containers.
 \end{otherlanguage*}
\end{resumo}